% OpenStax Calculus Volume 1 -- Section 4.5 Exercises: Derivatives and the Shape of a Graph
% Exercises reproduced from OpenStax, licensed under CC BY 4.0.
% Source: https://openstax.org/books/calculus-volume-1/pages/4-5-derivatives-and-the-shape-of-a-graph
\documentclass{exercises}
\title{OpenStax Calculus Volume 1}
\subtitle{Section 4.5 Exercises: Derivatives and the Shape of a Graph}
\sourceurl{https://openstax.org/books/calculus-volume-1/pages/4-5-derivatives-and-the-shape-of-a-graph}
\author{}
\date{}

\begin{document}
\maketitle

\begin{enumerate}[start=1]
\item If $c$ is a critical point of $f(x)$, when is there no local maximum or minimum at $c$? Explain.
\item For the function $y=x^{3}$, is $x=0$ both an inflection point and a local maximum/minimum?
\item For the function $y=x^{3}$, is $x=0$ an inflection point?
\item Is it possible for a point $c$ to be both an inflection point and a local extremum of a twice differentiable function?
\item Why do you need continuity for the first derivative test? Come up with an example.
\item Explain whether a concave-down function has to cross $y=0$ for some value of $x$.
\item Explain whether a polynomial of degree $2$ can have an inflection point.
\end{enumerate}

For the following exercises, analyze the graphs of $f'$, then list all intervals where $f$ is increasing or decreasing.

\begin{enumerate}[resume]
\item \textit{[Figure: The function $f'(x)$ is graphed. The function starts negative and crosses the $x$-axis at $(-2,0)$. Then it continues increasing a little before decreasing and crossing the $x$-axis at $(-1,0)$. It achieves a local minimum at $(1,-6)$ before increasing and crossing the $x$-axis at $(2,0)$.]}
\item \textit{[Figure: The function $f'(x)$ is graphed. The function starts negative and crosses the $x$-axis at $(-2,0)$. Then it continues increasing a little before decreasing and touching the $x$-axis at $(-1,0)$. It then increases a little before decreasing and crossing the $x$-axis at the origin. The function then decreases to a local minimum before increasing, crossing the $x$-axis at $(1,0)$, and continuing to increase.]}
\item \textit{[Figure: The function $f'(x)$ is graphed. The function starts negative and touches the $x$-axis at the origin. Then it decreases a little before increasing to cross the $x$-axis at $(1,0)$ and continuing to increase.]}
\item \textit{[Figure: The function $f'(x)$ is graphed. The function starts positive and decreases to touch the $x$-axis at $(-1,0)$. Then it increases to $(0,4.5)$ before decreasing to touch the $x$-axis at $(1,0)$. Then the function increases.]}
\item \textit{[Figure: The function $f'(x)$ is graphed. The function starts at $(-2,0)$, decreases to $(-1.5,-1.5)$, increases to $(-1,0)$, and continues increasing before decreasing to the origin. Then the other side is symmetric: that is, the function increases and then decreases to pass through $(1,0)$. It continues decreasing to $(1.5,-1.5)$, and then increase to $(2,0)$.]}
\end{enumerate}

For the following exercises, analyze the graphs of $f'$, then list all intervals where


(a) $f$ is increasing and decreasing and


(b) the minima and maxima are located.

\begin{enumerate}[resume]
\item \textit{[Figure: The function $f'(x)$ is graphed. The function starts at $(-2,0)$, decreases for a little and then increases to $(-1,0)$, continues increasing before decreasing to the origin, at which point it increases.]}
\item \textit{[Figure: The function $f'(x)$ is graphed. The function starts at $(-2,0)$, increases and then decreases to $(-1,0)$, decreases and then increases to an inflection point at the origin. Then the function increases and decreases to cross $(1,0)$. It continues decreasing and then increases to $(2,0)$.]}
\item \textit{[Figure: The function $f'(x)$ is graphed from $x=-2$ to $x=2$. It starts near zero at $x=-2$, but then increases rapidly and remains positive for the entire length of the graph.]}
\item \textit{[Figure: The function $f'(x)$ is graphed. The function starts negative and crosses the $x$-axis at the origin, which is an inflection point. Then it continues increasing.]}
\item \textit{[Figure: The function $f'(x)$ is graphed. The function starts negative and crosses the $x$-axis at $(-1,0)$. Then it continues increasing a little before decreasing and touching the $x$-axis at the origin. It increases again and then decreases to $(1,0)$. Then it increases.]}
\end{enumerate}

For the following exercises, analyze the graphs of $f'$, then list all inflection points and intervals where $f$ is concave up and concave down.

\begin{enumerate}[resume]
\item \textit{[Figure: The function $f'(x)$ is graphed. The function is linear and starts negative. It crosses the $x$-axis at the origin.]}
\item \textit{[Figure: The function $f'(x)$ is graphed. It is an upward-facing parabola with $0$ as its local minimum.]}
\item \textit{[Figure: The function $f'(x)$ is graphed. The function resembles the graph of $x^{3}$: that is, it starts negative and crosses the $x$-axis at the origin. Then it continues increasing.]}
\item \textit{[Figure: The function $f'(x)$ is graphed. The function starts negative and crosses the $x$-axis at $(-0.5,0)$. Then it continues increasing to $(0,1.5)$ before decreasing and touching the $x$-axis at $(1,0)$. It then increases.]}
\item \textit{[Figure: The function $f'(x)$ is graphed. The function starts negative and crosses the $x$-axis at $(-1,0)$. Then it continues increasing to a local maximum at $(0,1)$, at which point it decreases and touches the $x$-axis at $(1,0)$. It then increases.]}
\end{enumerate}

For the following exercises, draw a graph that satisfies the given specifications for the domain $x\in[-3,3]$. The function does not have to be continuous or differentiable.

\begin{enumerate}[resume]
\item $f(x)>0,f'(x)>0$ over $x>1,-3<x<0,f'(x)=0$ over $0<x<1$
\item $f'(x)>0$ over $x>2,-3<x<-1,f'(x)<0$ over $-1<x<2,f''(x)<0$ for all $x$
\item $f''(x)<0$ over $-1<x<1,f''(x)>0,-3<x<-1,1<x<3$, local maximum at $x=0$, local minima at $x=\pm2$
\item There is a local maximum at $x=2$, local minimum at $x=1$, and the graph is neither concave up nor concave down.
\item There are local maxima at $x=\pm1$, the function is concave up for all $x$, and the function remains positive for all $x$.
\end{enumerate}

For the following exercises, determine


(a) intervals where $f$ is increasing or decreasing and


(b) local minima and maxima of $f$.


If needed, use a calculator to graph the functions, but show your work.

\begin{enumerate}[resume]
\item $f(x)=\sin x+\sin^{3}x$ over $-\pi<x<\pi$
\item $f(x)=x^{2}+\cos x$
\end{enumerate}

For the following exercises, determine a. intervals where $f$ is concave up or concave down, and b. the inflection points of $f$.

\begin{enumerate}[resume]
\item $f(x)=x^{3}-4x^{2}+x+2$
\end{enumerate}

For the following exercises, determine


(a) intervals where $f$ is increasing or decreasing,


(b) local minima and maxima of $f$,


(c) intervals where $f$ is concave up and concave down, and


(d) the inflection points of $f$.

\begin{enumerate}[resume]
\item $f(x)=x^{2}-6x$
\item $f(x)=x^{3}-6x^{2}$
\item $f(x)=x^{4}-6x^{3}$
\item $f(x)=x^{11}-6x^{10}$
\item $f(x)=x+x^{2}-x^{3}$
\item $f(x)=x^{2}+x+1$
\item $f(x)=x^{3}+x^{4}$
\end{enumerate}

For the following exercises, determine


(a) intervals where $f$ is increasing or decreasing,


(b) local minima and maxima of $f$,


(c) intervals where $f$ is concave up and concave down, and


(d) the inflection points of $f$. Sketch the curve, then use a calculator to compare your answer. If you cannot determine the exact answer analytically, use a calculator.

\begin{enumerate}[resume]
\item \techrequired $f(x)=\sin(\pi x)-\cos(\pi x)$ over $x=[-1,1]$
\item \techrequired $f(x)=x+\sin(2x)$ over $x=[-\frac{\pi}{2},\frac{\pi}{2}]$
\item \techrequired $f(x)=\sin x+\tan x$ over $(-\frac{\pi}{2},\frac{\pi}{2})$
\item \techrequired $f(x)={(x-2)}^{2}{(x-4)}^{2}$
\item \techrequired $f(x)=\frac{1}{1-x},x\ne1$
\item \techrequired $f(x)=\frac{\sin x}{x}$ over $x=[-2\pi,0)\cup(0,2\pi]$
\item $f(x)=\sin(x)e^{x}$ over $x=[-\pi,\pi]$
\item $f(x)=(\ln x)\sqrt{x},x>0$
\item $f(x)=\frac{1}{4}\sqrt{x}+\frac{1}{x},x>0$
\item $f(x)=\frac{e^{x}}{x},x\ne0$
\end{enumerate}

For the following exercises, interpret the sentences in terms of $f$, $f'$, and $f''$.

\begin{enumerate}[resume]
\item The population is growing more slowly. Here $f$ is the population.
\item A bike accelerates faster, but a car goes faster. Here $f=$ Bike's position minus Car's position.
\item The airplane lands smoothly. Here $f$ is the plane's altitude.
\item A company tracks labor costs over the course of a year. On July 1, costs are at their peak. Here $f$ represents the cost of labor.
\item The economy is picking up speed. Here $f$ is a measure of the economy, such as GDP.
\end{enumerate}

For the following exercises, consider a third-degree polynomial $f(x)$, which has the properties $f'(1)=0,f'(3)=0$ (not applicable to Exercises 249 and 250). Determine whether the following statements are true or false. Justify your answer.

\begin{enumerate}[resume]
\item $f(x)=0$ for some $1\le x\le3$
\item $f''(x)=0$ for some $1\le x\le3$
\item There is no absolute maximum at $x=3$
\item If $f(x)$ has three roots, then it has $1$ inflection point.
\item If $f(x)$ has one inflection point, then it has three real roots.
\end{enumerate}

\end{document}
